\section{Auditoría de fuentes: convertir una entrevista con un CEO en un curso de Identity Systems}

Esta clase la imparte Todd McKinnon, cofundador y CEO de Okta. La sesión recorre la ventana para emprender en la computación en la nube de 2009, la infraestructura de identidad, los incidentes de seguridad, la cultura de la empresa, la adquisición de Auth0 y los AI agent. El repositorio conserva la portada oficial y 1,033 subtítulos con marca de tiempo. Este texto toma como eje de la clase la entrevista de febrero de 2025 y verifica los durable mechanism con materiales oficiales de Okta/Auth0 y con los estándares de NIST, IETF, OpenID y W3C. No presenta las capacidades de producto de 2026, como non-human identity, continuous session risk o Secure AI, como si ya estuvieran disponibles en el momento de la clase.

\begin{warningbox}{Límites de la evidencia sobre las cifras, los incidentes y los juicios organizacionales de la clase}
Las cifras de usuarios que inician sesión al mes, el volumen de solicitudes de autenticación, el número de empleados, la situación de los clientes y los juicios sobre el mercado que menciona el ponente se conservan \textbf{tal como se dijeron en clase} y solo sirven para ilustrar la escala de la workload y de la responsabilidad. La parte de incidentes de seguridad verifica los mecanismos de credential, session, logging y remediation con los RCA públicos, pero un incidente complejo no puede atribuirse a un solo empleado. Las opiniones sobre la adquisición y sobre la IA son juicios basados en la experiencia de McKinnon, no leyes causales universales.
\end{warningbox}

\begin{importantbox}{Los tres ciclos cerrados de esta clase}
\begin{enumerate}
  \item \textbf{Access loop}: identity source $\rightarrow$ policy $\rightarrow$ authentication/authorization $\rightarrow$ session/token $\rightarrow$ audit;
  \item \textbf{Trust loop}: incident $\rightarrow$ containment $\rightarrow$ root cause $\rightarrow$ remediation/notification $\rightarrow$ customer evidence;
  \item \textbf{Transition loop}: technology shift $\rightarrow$ focused wedge $\rightarrow$ infrastructure depth $\rightarrow$ organizational adaptation.
\end{enumerate}
\end{importantbox}

\begin{knowledgebox}{Términos clave: Identity no es una «tabla de cuentas»}
\term{principal} (sujeto) es una entidad que puede autenticarse y autorizarse; puede ser una persona, una service account, una workload o un agent. El \term{identity provider} (proveedor de identidad, IdP) mantiene el estado de los sujetos y proporciona a las aplicaciones aserciones de identidad confiables. La \term{authentication} (autenticación) responde «¿quién eres?»; la \term{authorization} (autorización) responde «¿qué puedes hacer con qué recurso?». La \term{session/token}, por su parte, convierte el resultado de una verificación en una capacidad de acceso reutilizable durante un tiempo limitado.
\end{knowledgebox}

\subsection{Por qué la página de inicio de sesión es solo la punta del iceberg}

Esta sección define primero el objeto de estudio del curso. Lo que el usuario ve es un login redirect o un MFA challenge; lo que la plataforma mantiene en realidad es directory, policy, protocol, session, provisioning, availability, security monitoring y audit. El servicio de identidad está en la «puerta de entrada» de muchas aplicaciones: un allow erróneo puede convertirse en una filtración de datos, y un deny erróneo o un outage puede impedir que toda la organización trabaje. Por eso pertenece a la vez a la infraestructura security-critical y availability-critical.

\subsection{Resumen de la sección}

El protagonista de la Lecture 09 no es la historia de emprendimiento de una empresa, sino el identity control plane. Todas las experiencias organizacionales que siguen se traducirán en preguntas de sistemas: dónde está el estado, cómo se toman las decisiones, cómo se acotan las fallas y cómo se recupera la confianza con evidencia.
